% TOP_PROBLEM: {"id":30007536,"problem_number":"LOCAL-30007536","title":"Sovereign–bank feedback loops","category_id":21,"category":{"id":21,"name":"economics","display_name":"Economics","description":"Open economic theory statements, published research agendas, and proposed scoped specifications; question provenance and historical priority are qualified per record.","slug":"economics","order_index":21,"created_at":"2026-10-08T00:00:00Z"},"status":"research_question","external_url":null,"legacy_ids":[],"published":false,"statement":"Find sovereign-exposure prudential policies that reduce endogenous sovereign–bank crisis losses while accounting for lending and safe-asset costs.","economics":{"original_id":25,"field":"Fiscal policy and sovereign debt","kind":"design","formalization_basis":"adapted_research_question","source_status":"research_agenda","priority_status":"earliest_found","priority_note":"This is the earliest explicit relevant statement verified in this audit, not a claim of original priority; older literature and earlier versions may contain antecedents.","earliest_verified_reference":{"authors":["Kris James Mitchener","Christoph Trebesch"],"title":"Sovereign Debt in the 21st Century","year":2021,"url":"https://www.ifo.de/DocDL/cesifo1_wp8959.pdf","doi":"10.3386/w28598","locator":"Section 10, “Concluding Remarks,” pp. 51–52","evidence_summary":"The survey says many puzzles remain, including how much more sovereign debt markets can absorb and how official holdings, bailouts, panics, and sovereign–bank feedback affect future debt crises."},"current_reference":{"authors":["Kris James Mitchener","Christoph Trebesch"],"title":"Sovereign Debt in the 21st Century","year":2021,"url":"https://www.ifo.de/DocDL/cesifo1_wp8959.pdf","doi":"10.3386/w28598","locator":"Section 10, “Concluding Remarks,” pp. 51–52","evidence_summary":"The survey says many puzzles remain, including how much more sovereign debt markets can absorb and how official holdings, bailouts, panics, and sovereign–bank feedback affect future debt crises."},"formalization_note":"A prudential-design problem narrows the survey’s explicitly continuing sovereign–bank crisis agenda; the precise policy vector and objective are newly specified."},"statusQualification":"Published question or research agenda verified at the cited locator. The mathematical specification is adapted for this catalogue and includes additional assumptions; source publication does not establish first-ever priority or certify this exact model remains unsolved. A prudential-design problem narrows the survey’s explicitly continuing sovereign–bank crisis agenda; the precise policy vector and objective are newly specified.","statusReviewedAt":"2026-10-08"}
% FIRST_POSTED: 2026-10-08
% ENTRY_KIND: statement_only
% !TeX program = lualatex
\documentclass[11pt]{article}
\usepackage{fontspec}
\usepackage[a4paper,margin=25mm]{geometry}
\usepackage{amsmath,amssymb,mathtools}
\usepackage{xurl}
\usepackage[hidelinks]{hyperref}
\newcommand{\sref}[2]{\hyperlink{source:#1:#2}{[S#2]}}
\setlength{\emergencystretch}{3em}
\begin{document}
% problemId:30007536
\section*{Sovereign–bank feedback loops}\label{30007536}
\subsection{Definitions and mathematical statement}
Consider a sovereign and competitive banks whose assets include loans \(L\) and domestic government bonds \(B^b\), funded by deposits and equity \(E\). Bank net worth changes with sovereign bond prices \(q^g\); loan supply depends on net worth and regulatory constraints. The government’s debt price in turn depends on fiscal capacity, output generated by lending, and possible bank support. Specify the joint default and bailout rules, household preferences, and the distribution of aggregate shocks before comparing policies.
Let a prudential policy \(a=(w,c,d)\) set the risk weight \(w\) on sovereign exposure, concentration limit \(c\), and required diversification \(d\). For feasible policies, define
\[
W(a)=\mathbb{E}\left[\sum_{t=0}^{H}\beta^t u(C_t(a),N_t(a))-\ell^{nc}(a)\right],
\]
where \(\ell^{nc}(a)\) denotes utility-equivalent nonconsumption resolution losses over the horizon, \(C_t\) is aggregate consumption after real resolution expenses, \(N_t\) labor, \(u\) household utility, and \(\beta\in(0,1)\) the discount factor. Nonconsumption resolution losses use the same utility scale and exclude expenses already deducted from consumption. Characterize policies that weaken the endogenous amplification from sovereign losses to banks and back without excessive reductions in productive credit or safe-asset services. Quantify the frontier between crisis-loss reduction and normal-time financing costs using identified exposure and lending responses. Address banks’ endogenous portfolio substitution and government incentives to provide guarantees. A policy comparison that treats sovereign bond risk and bank portfolios as fixed cannot establish its equilibrium welfare ranking.

\par\textbf{Formulation and scope.} A prudential-design problem narrows the survey’s explicitly continuing sovereign–bank crisis agenda; the precise policy vector and objective are newly specified.

\par\textbf{Open-question provenance.} An explicit question or research agenda was verified at the source locator. This does not by itself certify that every later version remains unresolved \sref{30007536}{1}.

\par\textbf{Historical priority.} This is the earliest explicit relevant statement verified in this audit, not a claim of original priority; older literature and earlier versions may contain antecedents.

\subsection{Short English statement}
Find sovereign-exposure prudential policies that reduce endogenous sovereign–bank crisis losses while accounting for lending and safe-asset costs.

\subsection{Sources}
\begin{itemize}\raggedright
\item[\textnormal{[S1]}]\hypertarget{source:30007536:1}{} Kris James Mitchener, Christoph Trebesch. 2021. Sovereign Debt in the 21st Century. Section 10, “Concluding Remarks,” pp. 51–52.
\url{https://www.ifo.de/DocDL/cesifo1_wp8959.pdf}.
DOI: 10.3386/w28598.
{\small\textit{Earliest verified question/agenda reference; historical priority qualified below. The survey says many puzzles remain, including how much more sovereign debt markets can absorb and how official holdings, bailouts, panics, and sovereign–bank feedback affect future debt crises.}}
\end{itemize}
\end{document}
